\section{Ordered transmission across a cut}\label{tr:section}

The pointwise comparison controls the effect of changing one small patch
of a prescribed boundary. We now turn this into a bound for an ordered
change along an entire separating cut. The quantity on the right will be
the change of its normalized magnetization. This lets a one-point
estimate control the law of every bounded observable beyond the cut.

\subsection{The cut geometry and the transmission estimate}

For the lattice statements, each local chart carries the critical
nearest-neighbor Ising model. No external fields are added beyond those
induced by prescribed spins, whose common signs remain arbitrary. We state the geometry
in a form which records all uniformity requirements.
A source $S$ and target $T$ are disjoint vertex sets. Common pins away
from them are allowed. Delete the common pinned vertices from the graph,
retaining $S$ as boundary vertices. A finite set $J$, disjoint from the
source, target, and common pins, is a \emph{separating cut} if every path
from $S$ to $T$ in this graph visits $J$. The continuation on the target
side of $J$ is then independent of the prescribed spins on $S$.

We will partition $J$ into small patches $b_i$. The pointwise comparison
bounds the effect of changing one patch; a straight counting segment
$L_i$ in the same neighborhood will record the corresponding magnetic
change. The next definition provides free space for this local count
and limits how often one spin can be counted.

\begin{definition}[An admissible cut at scales $w,s$]
\label{tr:admissible}
Fix constants $A,K\geq1$ and $a\in(0,1/(16A))$. For $0<w\leq s$, a separating cut
$J$ is admissible if it has a partition into patches $(b_i)$ and associated
straight counting segments $L_i\subseteq J$, with centers $z_i$ and
scales $\ell_i\in[aw,Aw]$, satisfying:
\begin{enumerate}
\item $b_i\cup L_i\subseteq z_i+B_{\ell_i}$ and
$|L_i|\in[a\ell_i,2\ell_i]$;
\item every $z_i+B_{16\ell_i}$ is a pin-free planar lattice chart,
including freedom from source pins;
\item each site of $J$ belongs to at most $K$ counting segments;
\item $T$ is at distance at least $s/A$ from every $z_i$,
$16\ell_i\leq s/A$, and the
patch comparison of Corollary~\ref{fd:patch-ratio} is applicable in a chart
up to a fixed positive multiple of $s$ at every patch.
\end{enumerate}
All constants in the following statements may depend on $a,A,K$ and the
fixed chart geometry, but not on $w,s$, the number of source pins, their
values, or the mesh. Decreasing $a$ or subdividing the patches by a fixed
factor preserves this convention.
\end{definition}

The principal examples are concentric square contours in a common free
collar of width comparable to $w$, separated from the target by a fixed
multiple of $s$, and parallel periodic rows with the same properties.
For a square, partition each side into intervals of length a sufficiently
small fixed multiple of $w$. Near a corner, choose a counting segment on
one of its adjacent sides. Taking these intervals smaller than the collar
width by a fixed factor gives the free squares in
Definition~\ref{tr:admissible}; a bounded number of neighboring intervals
can use the same segment, giving a uniform overlap bound. Fixed finite
unions of these cuts are allowed. The hypotheses are also appropriate
when $w/s$ is small: the constants just described do not change with that
ratio.

Write
\begin{equation}
\label{tr:magnetization-definition}
 M_J=s^{-7/8}\sum_{x\in J}\sigma_x.
\end{equation}
The normalization always uses the displayed reference scale $s$, not
the individual patch scale.

\begin{proposition}[Transmission]
\label{tr:transmission}
There is an exponent $\delta\in(0,1/8)$, chosen by decreasing the exponent
of Corollary~\ref{fd:patch-ratio} if necessary, with the following property.
Prescribe two ordered configurations on $S$, use identical common pins
elsewhere, and let $\Delta$ denote higher-law expectation minus lower-law
expectation. For an admissible separating cut at scales $w,s$, a bounded
observable $h$ supported on $T$, and $w$ sufficiently large,
\begin{equation}
\label{tr:transmission-eq}
 |\Delta h|
 \leq C\|h\|_\infty (w/s)^{-7/8+\delta}\Delta M_J.
\end{equation}
Equivalently, if $\|h\|_\infty\leq H$ for fixed $H>0$,
\begin{equation}
\label{tr:transmission-lower}
 \Delta M_J\geq c_H (w/s)^{7/8-\delta}|\Delta h|.
\end{equation}
The constants are uniform as $w/s$ decreases within the stated geometry.
\end{proposition}

The exponent in \eqref{tr:transmission-eq} comes from two estimates.
A changed patch of diameter comparable to $w$ changes the conditional
expectation of $h$ by at most $C\|h\|_\infty(w/s)^\delta$, by the
small-patch comparison. We construct an ordered coupling in which,
conditional on a disagreement reaching that patch, the expected number
of disagreements on its counting segment is at least $cw^{7/8}$. Bounded overlap sums these contributions over
$J$. The normalization $s^{-7/8}$ then gives the factor
$(w/s)^{-7/8+\delta}$.

We prove the second estimate first. An auxiliary FK representation
will identify ordered disagreements with connections to the changed
pins. A local domination statement will let us test for magnetic mass
inside the pin-free charts of Definition~\ref{tr:admissible}. We will
then return to the cut and prove the displayed transmission estimate.
\subsection{An auxiliary FK model for ordered disagreements}
\label{tr:doubled-subsection}

Begin on a finite graph $\mathcal G=(V,E)$ with couplings $K_e\geq0$
and a pin set $P\subseteq V$. For a retained free set $H\subseteq V\setminus P$,
at specified couplings $J_e$ let $Z_{P,\tau}(H)$ denote the Ising
partition sum on the induced graph $P\cup H$, with signs $\tau$ on $P$. Partition sums include edges between
pinned vertices. Their contribution is constant in the free spins.
In the planar applications the original couplings are $K_e=\beta_c$.
The FK bonds introduced below instead have doubled couplings on a
random induced set; the critical FK crossing theorem is not applied
directly to these auxiliary bonds.

Take two independent Ising configurations $\eta,\xi$ whose signs on the
same pin set $P$ are ordered, with the higher signs assigned to $\eta$.
Let $U\subseteq P$ be the disagreement pins and $A=P\setminus U$ the
agreement pins, whose common signs are denoted by $\tau$. Let $W=V\setminus
P$, and let $G\subseteq W$ be the set of free sites where the original
replicas disagree. Thus the full auxiliary disagreement set is $G\cup U$.

\begin{lemma}[Doubled representation and ordered coupling]
\label{tr:doubled-coupling}
With partition sums at couplings $J_e=2K_e$, the law of $G$ is
\begin{equation}
\label{tr:doubled-weight}
 \mu_\tau(G)\ \propto\ Z_{U,+}(G)\,Z_{A,\tau}(W\setminus G).
\end{equation}
Conditional on $G$, the spins $\eta=-\xi$ on $G\cup U$ form the first
Ising law in \eqref{tr:doubled-weight}, independently of the common spins
on its complement. Put FK bonds $\omega$ on $G\cup U$. There is a coupling
of the original two Ising marginals, denoted again by $(\eta,\xi)$, for
which $\eta\geq\xi$ and
\begin{equation}
\label{tr:mean-connection}
 \eta_x-\xi_x=2\mathbf1_{\{x\longleftrightarrow U\}},
 \qquad
 \mathbb E\eta_x-\mathbb E\xi_x
 =2\mathbb P(x\longleftrightarrow U).
\end{equation}
Here and below a connection to $U$ uses actual open edges, and is false for
$x\notin G\cup U$.
\end{lemma}

\begin{proof}
For an edge with endpoints in the same one of $G\cup U$ and its complement,
the two replica energies add to twice the energy of their common
orientation. For an edge between these sets they cancel. This proves
\eqref{tr:doubled-weight} and the conditional independence. Edges with
both endpoints pinned contribute constants independent of $G$, so cause
no exception to the factorization.

Every FK cluster in $G\cup U$ avoiding $U$ has an independent fair
orientation. On each such cluster replace the pair of opposite signs by
two copies of the first replica's sign. The first replica is unchanged;
the second replica's collection of independent fair cluster signs has been
negated, so its conditional marginal is unchanged as well. Clusters
meeting $U$ retain the pair $(+,-)$, and agreement sites retain equal
signs. This proves both the ordering and \eqref{tr:mean-connection}.
The auxiliary set $G$ need not equal the disagreement set after this
operation: the latter is precisely the union of clusters meeting $U$.
\end{proof}

Order pairs $(G,\omega)$ by inclusion in both coordinates, padding the
bond configuration by closed edges off $G\cup U$. An event or function of
this pair is called increasing in this order. The \emph{plus-agreement
reference} is the law in which all signs on $A$ have been changed to plus,
while the disagreement pins $U$ are retained.

The coupled replicas now differ exactly on clusters meeting $U$.
The auxiliary set $G$ still records the original independent replicas,
and can contain clusters on which the coupled replicas agree. To obtain
lower bounds on its local connections, we compare its law with the
plus-agreement reference. The density between the two auxiliary set laws
is, up to normalization,
\[
 r(G)=\frac{Z_{A,\tau}(W\setminus G)}{Z_{A,+}(W\setminus G)}.
\]
We will show that this density increases when $G$ grows. The needed
partition-sum fact is a comparison of the factors gained by adding one
free vertex. We prove it for arbitrary nonnegative couplings $J_e$,
then apply it with $J_e=2K_e$.

\begin{lemma}[Spin products and addition of vertices]
\label{tr:site-gain}
Fix a set $P$ of pinned vertices, couplings $J_e\geq0$, and let
$Z_{P,\tau}(H)$ be the Ising partition sum on the induced graph with vertex
set $P\cup H$, with signs $\tau$ on $P$. Here $H\subseteq V\setminus P$ is
the retained free set. Then the following assertions hold.
\begin{enumerate}
\item With all pins plus, every spin-product expectation is nonnegative and
increases when nonnegative couplings are increased, free vertices and their
incident ferromagnetic edges are added, or additional spins are pinned plus.
The product in question involves only vertices present before the change.
\item For every finite set $Q\subseteq P\cup H$,
\begin{equation}
\label{tr:mixed-product}
 \left|\mathbb E_{P,\tau,H}\prod_{x\in Q}\sigma_x\right|
 \leq \mathbb E_{P,+,H}\prod_{x\in Q}\sigma_x.
\end{equation}
\item The function $H\mapsto Z_{P,+}(H)$ is log-supermodular. Moreover,
for $v\notin P\cup H$, the site gains satisfy
\begin{equation}
\label{tr:gain-comparison}
 \frac{Z_{P,\tau}(H\cup\{v\})}{Z_{P,\tau}(H)}
 \leq
 \frac{Z_{P,+}(H\cup\{v\})}{Z_{P,+}(H)}.
\end{equation}
\end{enumerate}
\end{lemma}

\begin{proof}
Use the FK representation with edge parameters $p_e=1-e^{-2J_e}$. In
descriptions of sign compatibility, an \emph{actual cluster} means a
component formed by open edges of the original graph: virtual wires between
pins do not count as open paths. Write $k_0(\omega)$ for the number of actual
clusters avoiding $P$. Apart from a constant, the all-plus bond weight is
\begin{equation}
\label{tr:pinned-fk-weight}
 \prod_{e\in E(P\cup H)}p_e^{\omega_e}(1-p_e)^{1-\omega_e}
 2^{k_0(\omega)}.
\end{equation}
This is the usual FK law with all pins wired. Conditional on the bonds,
clusters touching a pin have sign plus, and the other clusters have
independent fair signs.

Let $\mathcal F_Q$ be the event that each actual cluster avoiding $P$
contains an even number of vertices of $Q$. Conditional sign averaging gives
\begin{equation}
\label{tr:parity-event}
 \mathbb E_{P,+,H}\prod_{x\in Q}\sigma_x
 =\phi_{P,+,H}(\mathcal F_Q).
\end{equation}
The event $\mathcal F_Q$ is increasing. Indeed, merging two unpinned
clusters whose parities are even preserves their even parity, whereas
merging any cluster with a pinned cluster removes its parity restriction.
Adding an initially isolated vertex outside $Q$ adds a cluster of even
parity zero. Thus FK monotonicity under edge addition, increases of edge
parameters, and increases of the boundary wiring proves the first
assertion. Pinning an additional spin plus is covered by the same argument:
its cluster ceases to have a parity restriction and the FK wiring increases.
Positive fields, when needed, may equivalently be represented by edges to
one additional plus vertex.

For mixed signs, put
\[
 \mathcal D_\tau=
 \{\text{no actual cluster contains two pins with opposite signs}\}.
\]
The mixed-sign bond weight is exactly
\eqref{tr:pinned-fk-weight} multiplied by $\mathbf1_{\mathcal D_\tau}$.
In particular, its law is $\phi_{P,+,H}(\,\cdot\mid\mathcal D_\tau)$.
This identity also handles edges between opposite pins: such an edge must
be closed. The event $\mathcal D_\tau$ is decreasing and has positive
probability. Conditional sign averaging in the mixed law has absolute
value at most $\mathbf1_{\mathcal F_Q}$. Positive association of the wired
FK law therefore gives
\[
 \left|\mathbb E_{P,\tau,H}\prod_{x\in Q}\sigma_x\right|
 \leq \phi_{P,+,H}(\mathcal F_Q\mid\mathcal D_\tau)
 \leq \phi_{P,+,H}(\mathcal F_Q),
\]
which proves \eqref{tr:mixed-product}.

For the last assertion, summing the spin at a newly added vertex gives the
exact identity
\begin{equation}
\label{tr:site-gain-formula}
 g_{P,\tau,H}(v):=
 \frac{Z_{P,\tau}(H\cup\{v\})}{Z_{P,\tau}(H)}
 =2\mathbb E_{P,\tau,H}
 \cosh\!\left(\sum_{u\in P\cup H}J_{vu}\sigma_u\right).
\end{equation}
Zero couplings may be omitted from the sum. If $N$ is its neighbor set,
the integrand has the expansion
\begin{equation}
\label{tr:cosh-expansion}
 \cosh\!\left(\sum_{u\in N}J_{vu}\sigma_u\right)
 =\prod_{u\in N}\cosh J_{vu}
 \sum_{\substack{Q\subseteq N\\ |Q|\ {\rm even}}}
 \prod_{u\in Q}\tanh J_{vu}\prod_{u\in Q}\sigma_u.
\end{equation}
All coefficients are nonnegative. Under plus pins, enlarging $H$ increases
the old product expectations, adds nonnegative terms, and increases the
prefactor. Hence $g_{P,+,H}(v)$ increases with $H$. Increasing site gains are
equivalent to log-supermodularity: add the vertices of $H_1\setminus H_2$
one at a time, comparing their gains above $H_1\cap H_2$ and above $H_2$,
to obtain
\[
 Z_{P,+}(H_1\cup H_2)Z_{P,+}(H_1\cap H_2)
 \geq Z_{P,+}(H_1)Z_{P,+}(H_2).
\]
Finally, apply \eqref{tr:mixed-product} term by term in
\eqref{tr:cosh-expansion} to prove \eqref{tr:gain-comparison}.
\end{proof}

\begin{lemma}[Association and increasing conditioning]
\label{tr:increasing-conditioning}
The plus-agreement reference law of $(G,\omega)$ is positively associated.
The mixed-agreement law has an increasing density, depending only on $G$,
relative to this reference. Consequently, if $\mathcal E$ is any increasing
event of positive probability and $f$ is any increasing function, then
\begin{equation}
\label{tr:conditional-reference}
 \mathbb E_\tau[f\mid\mathcal E]\geq\mathbb E_+[f].
\end{equation}
There is no restriction on the number or placement of the pins in this
finite-graph assertion.
\end{lemma}

\begin{proof}
By Lemma~\ref{tr:site-gain}, both factors in
\eqref{tr:doubled-weight} with plus pins are log-supermodular functions of
$G$. Complementation exchanges unions and intersections and hence
preserves this property. Their product is log-supermodular, so the
reference set law is positively associated. Conditional on $G$, the FK
law is positively associated and increases stochastically with $G$.
Indeed, one may first add isolated vertices and then increase their edge
parameters from zero. Thus conditional expectations of increasing
functions of $(G,\omega)$ increase with $G$. The decomposition of a
covariance into its conditional covariance and the covariance of its
conditional expectations proves joint positive association.

The mixed law has the density $r(G)$ displayed above.
By \eqref{tr:gain-comparison}, the quotient $Z_{A,\tau}(H)/Z_{A,+}(H)$
decreases as $H$ grows. Therefore $r$ increases with $G$. The conditional
law on the left side of \eqref{tr:conditional-reference} has density
proportional to $r\mathbf1_{\mathcal E}$ against the associated reference.
That density is increasing, and association gives
\[
 \frac{\mathbb E_+[f r\mathbf1_{\mathcal E}]}
 {\mathbb E_+[r\mathbf1_{\mathcal E}]}
 \geq\mathbb E_+[f].
\]
This argument requires neither association of the mixed law nor
preservation of a lattice condition after increasing-event conditioning.
\end{proof}

\begin{proposition}[Local domination]
\label{tr:local-domination}
Let $Q\subseteq W$ contain no pins. In $Q$, form the doubled model from two
independent Ising configurations with plus signs on every exterior
neighbor of $Q$, with the original couplings $K_e$. Call its auxiliary
set-and-bond law $\mathbb P_Q^{++}$. For every increasing function $f$ of
the auxiliary set and bonds internal to $Q$, and every positive-probability
increasing event $\mathcal E$ in the original doubled model,
\begin{equation}
\label{tr:local-domination-eq}
 \mathbb E_\tau[f\mid\mathcal E]\geq\mathbb E_Q^{++}[f].
\end{equation}
The same statement holds with $\mathcal E$ the sure event.
\end{proposition}

\begin{proof}
By Lemma~\ref{tr:increasing-conditioning}, it suffices to compare the
plus-agreement reference with $\mathbb P_Q^{++}$. Condition its exterior
auxiliary set $G\setminus Q$. Fix $v\in Q$ and the membership of all other
sites of $Q$. Write $D$ for the full disagreement vertices excluding $v$,
including the pins $U$, and $H$ for the full agreement vertices excluding
$v$, including the pins $A$. The conditional odds for membership of $v$
are the quotient of two site gains:
\begin{equation}
\label{tr:local-odds}
 \frac{\mu_+(v\in G\mid G\setminus\{v\})}
 {\mu_+(v\notin G\mid G\setminus\{v\})}
 =\frac{g_{U,+,D\setminus U}(v)}
             {g_{A,+,H\setminus A}(v)}.
\end{equation}
Here $g$ is the gain in \eqref{tr:site-gain-formula}.

The numerator is at least the gain on the induced disagreement set in
$Q$: adding the exterior disagreement vertices, their edges, and their
plus pins only increases it. For the denominator, first pin every exterior
agreement vertex plus. The parity argument in Lemma~\ref{tr:site-gain}
and the nonnegative expansion \eqref{tr:cosh-expansion} show that this
increases the gain. Then restore every missing exterior neighbor of $Q$
as a plus pin, again increasing the gain. The resulting denominator is
exactly the agreement-site gain for the box-alone model: partition
factors involving only the now-pinned exterior cancel in the site ratio.
Thus the odds
in \eqref{tr:local-odds} are at least the corresponding odds of
$\mathbb P_Q^{++}$ at every fixed interior set configuration.

Equivalently, the density of the reference interior-set law relative to
the box-alone set law is increasing: the ratio increases under every
single-site addition. The box-alone set law is associated by
Lemma~\ref{tr:increasing-conditioning}. Tilting it by this increasing
density therefore proves stochastic domination of the sets.

Given a full auxiliary set, conditioning the exterior FK bonds induces
some wiring on its part in $Q$. This wiring dominates the free boundary
condition of the box-alone disagreement-bond system; its orientation has
no disagreement pins in $Q$ or on its exterior. The free FK law on an
induced set also grows with that set. First perform this bond comparison
and then the set comparison just obtained. Averaging over the exterior
set proves the claim for every increasing $f$.
\end{proof}

\subsection{Mass carried by a disagreement connection}
\label{tr:mass-subsection}

Local domination remains valid after conditioning on a connection from
the disagreement pins. We now construct, in the doubled plus-boundary
reference, an increasing event joining many sites of a segment to such
a connection. This is where the correlation and moment estimates of
Section~\ref{mo:section} enter.

For the remainder of the section the underlying graph in each local
chart is the square lattice at $\beta_c$. Bounded integer roundings of
rectangles and segments are understood. Every lower bound below holds
once its scale is sufficiently large; the threshold depends only on the
fixed geometric constants. This is sufficient when macroscopic geometry
is fixed before taking the mesh limit.

The bonds in this subsection belong to the auxiliary set sampled before
the recoloring in Lemma~\ref{tr:doubled-coupling}. Equal plus boundary
assignments can therefore yield auxiliary bond connections, even though
the resulting ordered spin coupling has identical replicas.

\begin{lemma}[Many sites joined across a rectangle]
\label{tr:many-sites}
Consider two independent plus-boundary Ising configurations in a rectangle
of scale $w$ and fixed aspect ratio, and their auxiliary disagreement
bonds. Let $L_1,L_2$ be straight lattice segments of lengths comparable to
$w$, in fixed compact bulk subsets of the rectangle, mutually separated
by a distance comparable to $w$. There are constants $c,p>0$, depending
only on this geometry, such that
\begin{equation}
\label{tr:many-sites-eq}
 \mathbb P_Q^{++}\bigl(
 |\{x\in L_1:x\longleftrightarrow L_2\}|\geq c w^{7/8}
 \bigr)\geq p.
\end{equation}
In particular a disagreement-bond connection between the segments has
probability bounded below.
\end{lemma}

\begin{proof}
In the original independent pair put $Y_x=(\eta_x-\xi_x)/2$. There are no
disagreement pins. Conditional cluster signs give
\begin{equation}
\label{tr:pair-identity}
 \mathbb P_Q^{++}(x\longleftrightarrow y)
 =\mathbb E[Y_xY_y]
 =\tfrac12\operatorname{Cov}_Q^+(\sigma_x,\sigma_y).
\end{equation}
More generally,
\begin{equation}
\label{tr:four-bound}
 \mathbb P_Q^{++}(x\longleftrightarrow y,
                  x'\longleftrightarrow y')
 \leq\mathbb E[Y_xY_yY_{x'}Y_{y'}].
\end{equation}
Indeed, if both pairs connect, all auxiliary clusters contain an even
number of these four indexed factors. Conditional sign averaging of the
product is the indicator of this even-parity event, with the additional
requirement that all factors lie in the auxiliary set. This explanation
also proves \eqref{tr:four-bound} when indices repeat.

Set
\[
 X=\sum_{x\in L_1}\sum_{y\in L_2}
        \mathbf1_{\{x\longleftrightarrow y\}}.
\]
The bulk connected-correlation lower bound in
Lemma~\ref{fd:static-correlations} gives $\mathbb EX\geq c_1w^{7/4}$.
By \eqref{tr:four-bound}, Cauchy--Schwarz, and the fourth line moments in
Lemma~\ref{fd:moments},
\[
 \mathbb EX^2
 \leq\mathbb E\left[
       \left(\sum_{L_1}Y_x\right)^2
       \left(\sum_{L_2}Y_y\right)^2\right]
 \leq C_1 w^{7/2}.
\]
The moments for $Y$ follow from those for either replica by the triangle
inequality in $L^4$. Paley--Zygmund therefore gives constants $c_2,p_2>0$
such that $\mathbb P(X\geq c_2w^{7/4})\geq p_2$.

Let $L$ be the largest number of sites of $L_2$ in one auxiliary cluster.
Then
\[
 L^2\leq\sum_{\mathcal C}|\mathcal C\cap L_2|^2,
 \qquad
 \mathbb E\sum_{\mathcal C}|\mathcal C\cap L_2|^2
 =\mathbb E\left(\sum_{L_2}Y_x\right)^2
 \leq C_2w^{7/4}.
\]
Choose $K$ so large that $C_2/K^2<p_2/2$. With probability at least
$p_2/2$, both $X\geq c_2w^{7/4}$ and $L\leq Kw^{7/8}$ hold. Every site
of $L_1$ reaching $L_2$ contributes at most $L$ to $X$. On that
intersection at least $(c_2/K)w^{7/8}$ sites reach $L_2$, proving
\eqref{tr:many-sites-eq}. Although the event involving $L$ was used to
establish its probability, the event in \eqref{tr:many-sites-eq} itself
is increasing.
\end{proof}

\begin{lemma}[A surrounding disagreement circuit]
\label{tr:circuit}
In the doubled plus-boundary model on $B_{16w}$, a circuit of actual
disagreement bonds in $B_{4w}\setminus B_{3w/2}$ surrounding $B_{3w/2}$ has
probability at least a constant $p_0>0$ for all sufficiently large $w$.
\end{lemma}

\begin{proof}
Consider the four strips
\[
 [-4w,4w]\times[2w,4w],\quad
 [-4w,4w]\times[-4w,-2w],\quad
 [2w,4w]\times[-4w,4w],\quad
 [-4w,-2w]\times[-4w,4w].
\]
For the top strip, extend the long coordinate to $[-5w,5w]$ and put
agreeing plus boundary conditions on this extended rectangle. Apply
Lemma~\ref{tr:many-sites} to vertical segments at $x=-4w$ and $x=4w$,
each with its $y$ coordinate in $[5w/2,7w/2]$. The segments are in its
bulk. Their connection stays inside the extended rectangle, since the
auxiliary disagreement set cannot contain its agreeing boundary pins.
A subpath of this connection gives a left-to-right crossing of the
unextended top strip. Proposition~\ref{tr:local-domination} transfers
its positive lower bound to the doubled model on $B_{16w}$. The same
construction applies to the other three strips.

The four crossing events are increasing, so their intersection has
probability at least the product of their four lower bounds. On this
intersection, adjacent crossings intersect inside their corner overlap:
one crosses that overlap in one coordinate and the other in the
perpendicular coordinate. The union contains a circuit surrounding the
central square. One way to see the last assertion is to join successive
intersection points along the four crossings, erase portions lying beyond
the adjacent corner overlaps, and obtain a closed walk in the annulus
with winding number one about the origin. Removing closed subwalks
decomposes it into simple circuits; at least one retains nonzero winding
number. All its edges lie in the four strips. This proves the claim.
\end{proof}

\begin{proposition}[Mass at a source-connected patch]
\label{tr:mass}
Fix $a>0$. Let $b_0\subseteq z+B_w$, and let $L\subseteq z+B_w$ be a
horizontal or vertical lattice segment of length at least $aw$ and at
most $2w$. Suppose $z+B_{16w}$ is a free square-lattice neighborhood:
it contains no common or disagreement pins. Then in the doubled model
with arbitrary ordered pins outside this square,
\begin{equation}
\label{tr:mass-eq}
 \sum_{x\in L}\mathbb P(x\longleftrightarrow U)
 \geq c_a w^{7/8}\mathbb P(U\longleftrightarrow b_0).
\end{equation}
The constant and lower threshold for $w$ depend only on $a$. The same
statement holds inside an isometric planar chart of a torus or cylinder.
\end{proposition}

\begin{proof}
Translate the picture so that $z=0$. Choose a translate $L'$ of $L$
by $6w$ in a perpendicular coordinate. Both segments lie in fixed bulk
subsets of $B_{16w}$, while $L'$ lies outside $B_{4w}$. Under its doubled
plus-boundary law, Lemma~\ref{tr:many-sites} gives positive probability
that at least $c w^{7/8}$ sites of $L$ reach $L'$ by paths using bonds
internal to $B_{16w}$. Intersect this increasing local event with the
increasing circuit event of Lemma~\ref{tr:circuit}.
Association makes the probability of their intersection at least a
positive constant $p$.

If $\mathbb P(U\longleftrightarrow b_0)=0$, the claim is immediate.
Otherwise condition on that increasing event and use
Proposition~\ref{tr:local-domination}. The circuit and the many-sites
event still occur jointly with conditional probability at least $p$.
Every path from $U$, which is outside $B_{16w}$, to $b_0\subseteq B_w$
meets the surrounding circuit. Every counted path from $L\subseteq B_w$
to $L'$ also meets it. By planarity all the counted sites are therefore
connected to $U$. The conditional expected number of such sites is at
least $pc w^{7/8}$. Multiplying by the probability of the conditioning
event proves \eqref{tr:mass-eq}.
\end{proof}

\subsection{From local mass to transmission}
\label{tr:transmission-subsection}

Proposition~\ref{tr:mass} provides the local estimate promised at the
start of the section. At each patch of an admissible cut, its scale
$\ell_i\asymp w$ converts the probability of reaching that patch into
expected disagreement mass on $L_i$. We now sum these estimates and
use the common target-side continuation to complete the proof.

\begin{proof}[Proof of Proposition~\ref{tr:transmission}]
Use the monotone coupling of Lemma~\ref{tr:doubled-coupling}. Let
$\eta_J\geq\xi_J$ be its cut configurations. Put
$k(j)=\mathbb E[h\mid\sigma_J=j]$, where the conditional expectation
uses the common target-side continuation. The spatial Markov property
and the separating property imply that this function does not depend
on which source configuration was prescribed. Consequently
\[
 \Delta h=\mathbb E[k(\eta_J)-k(\xi_J)].
\]

In this telescope, the other cut spins are common pins for the pointwise
comparison. The local mass estimate is applied to the original coupling,
whose counting charts remain pin-free.

Interpolate between the two cut configurations one patch at a time,
using any deterministic ordering of their partition. Intermediate
configurations are legitimate cut configurations; spins at all other
patches are fixed for each comparison. Corollary~\ref{fd:patch-ratio}
bounds the oscillation of $k$ in one patch by
$C\|h\|_\infty(w/s)^\delta$. It applies with these arbitrary other cut
spins as common pins. If the two sampled cut configurations agree on
the patch, the corresponding summand is zero. By
\eqref{tr:mean-connection}, disagreement occurs somewhere in $b_i$
exactly when $U\longleftrightarrow b_i$. Therefore
\begin{equation}
\label{tr:telescope}
 |\Delta h|\leq C\|h\|_\infty(w/s)^\delta
 \sum_i\mathbb P(U\longleftrightarrow b_i).
\end{equation}

Proposition~\ref{tr:mass}, applied with scale $\ell_i\asymp w$, and the
overlap bound on the counting segments give
\[
 \sum_i\mathbb P(U\longleftrightarrow b_i)
 \leq C w^{-7/8}\sum_i\sum_{x\in L_i}
               \mathbb P(U\longleftrightarrow x)
 \leq C K w^{-7/8}\sum_{x\in J}
               \mathbb P(U\longleftrightarrow x).
\]
The last sum is $\tfrac12s^{7/8}\Delta M_J$ by
\eqref{tr:mean-connection}. Substitute it in \eqref{tr:telescope} to
obtain \eqref{tr:transmission-eq}. All local constructions use the same
fixed aspect ratios and overlap constants, independently of $w/s$.
\end{proof}

\begin{definition}[Order-positive signed source]
\label{tr:order-positive}
For a probability law $\mu$ on the source spins, a real function $q$ is
an order-positive centered density if $\mathbb E_\mu q=0$ and
$\mathbb E_\mu[qf]\geq0$ for every increasing source function $f$.
The word density here refers to a signed measure; $q$ need not be
nonnegative or itself increasing.
\end{definition}

\begin{corollary}[Random and signed sources]
\label{tr:signed-transmission}
The conclusions of Proposition~\ref{tr:transmission} hold after averaging
over any ordered coupling of two random source configurations, with the
appropriate common Gibbs continuation given each source. In particular,
let $\mu$ be a Gibbs law with common pins and let $q=q(\sigma_S)$ be an
order-positive centered density for its source marginal. Then
\begin{equation}
\label{tr:signed-transmission-eq}
 \left|\mathbb E_\mu[qh]\right|
 \leq C\|h\|_\infty(w/s)^{-7/8+\delta}
                    \mathbb E_\mu[qM_J].
\end{equation}
The constants do not depend on the mass of the positive part of $q$.
\end{corollary}

\begin{proof}
For an ordered random pair of sources, apply
\eqref{tr:transmission-eq} to each pair and integrate. The magnetic
differences are nonnegative. The triangle inequality bounds the absolute
value of the integrated target difference by the integral of its absolute
value, proving the first assertion.

For the signed assertion put $m=\mathbb E_\mu q_+=\mathbb E_\mu q_-$.
If $m=0$ there is nothing to prove. Otherwise let
$\mu_\pm=m^{-1}q_\pm\mu_S$ be probability laws on source configurations.
Order positivity says $\mu_+(F)\geq\mu_-(F)$ for every increasing set $F$.
These laws admit a coupling $(s_+,s_-)$ with $s_+\geq s_-$. For
completeness, this finite transport fact follows by sending mass
$\mu_-(x)$ from each lower configuration $x$ to upper configurations
$y\geq x$ with capacities $\mu_+(y)$. For a set $B$ of lower
configurations, the required cut inequality is
\[
 \mu_-(B)\leq\mu_-(\mathord\uparrow B)
             \leq\mu_+(\mathord\uparrow B),
\]
so the finite max-flow/min-cut criterion transports all the mass.
Apply the random-source assertion to this coupling and multiply by $m$.
The Gibbs continuation changes the resulting difference of source
integrals into $\mathbb E_\mu[qh]$ and $\mathbb E_\mu[qM_J]$ respectively,
which proves \eqref{tr:signed-transmission-eq}.
\end{proof}

\begin{corollary}[Ordered changes of domains and total variation]
\label{tr:ordered-tv}
Suppose two ordered Ising laws, possibly defined by different sets of
pins or different volumes, have a common free source contour $S_0$ such
that their conditional continuation from $S_0$ toward a target $T$ is the
same. Let an admissible cut $J$ lie between $S_0$ and $T$. Then
\eqref{tr:transmission-eq} holds for their expectation differences.
In particular, with total variation defined as the supremum of event
probability differences,
\begin{equation}
\label{tr:ordered-tv-eq}
 \|\mu_T^+-\mu_T^-\|_{\rm TV}
 \leq C(w/s)^{-7/8+\delta}
       \bigl(\mathbb E_{\mu^+}M_J-\mathbb E_{\mu^-}M_J\bigr).
\end{equation}
At fixed positive $w/s$, convergence of the magnetic difference to zero
therefore implies convergence in total variation on the buffered target.
\end{corollary}

\begin{proof}
Restrict an ordered coupling to $S_0$ and then use the common Gibbs
continuation. Corollary~\ref{tr:signed-transmission} in its random-source
form applies. Taking $h$ to be an arbitrary target-event indicator proves
\eqref{tr:ordered-tv-eq}. All differing pins must be on the source side of
$S_0$; the common-continuation hypothesis explicitly excludes a changed
pin in the target component.
\end{proof}

